% =====================================================================
\chapter{Forma polar, fórmula de Euler, De Moivre y rotaciones}
\label{tema:6}
% =====================================================================

La forma binómica es cómoda para sumar; la forma polar, para multiplicar. Este tema demuestra que al multiplicar dos complejos se multiplican los módulos y se suman los argumentos, introduce la exponencial imaginaria $e^{i\theta}$ y la fórmula de Euler, deduce el teorema de De Moivre y la fórmula de las raíces $n$-ésimas, y termina interpretando la multiplicación por $e^{i\alpha}$ como el giro de ángulo $\alpha$ del plano, que es la idea sobre la que se construye RoPE. Se siguen la sección~6 de las notas de 18.03 del MIT~\cite{mit1803}, las §1.2 a §1.5 de las notas de Cambridge~\cite{cambvm}, la §10.5 de~\cite{osat} y la lectura geométrica del capítulo~1 de Needham~\cite{needham}.

\section{Argumento y forma polar}

\begin{definicion}[Argumento y forma polar]\label{def:polar}
Sea $z\ne0$ y $r=\abs{z}$. Un \emph{argumento} de $z$ es cualquier $\theta\in\R$ tal que
\[
z=r(\cos\theta+i\sin\theta).
\]
Esta expresión es la \emph{forma polar} de $z$. El \emph{argumento principal} $\Arg z$ es el único argumento en $(-\pi,\pi]$. El número $0$ no tiene argumento.
\end{definicion}

Todo $z\ne0$ tiene argumento: $z/r$ tiene módulo~1, luego es un punto $P(\theta)$ de la circunferencia unidad. Por el \cref{teo:simetrias}(ii), dos argumentos de $z$ difieren en un múltiplo entero de $2\pi$, y exactamente uno está en $(-\pi,\pi]$.

\begin{proposicion}[Cálculo del argumento principal]\label{prop:arg}
Sea $z=a+bi\ne0$. Si $a>0$, $\Arg z=\arctan(b/a)$. Si $a<0$, $\Arg z=\arctan(b/a)+\pi$ cuando $b\ge0$ y $\arctan(b/a)-\pi$ cuando $b<0$. Si $a=0$, $\Arg z=\pi/2$ cuando $b>0$ y $-\pi/2$ cuando $b<0$.
\end{proposicion}

\addtocounter{definicion}{1}%

\section{El producto en forma polar}

\begin{teorema}[Producto y cociente en forma polar]\label{teo:producto-polar}
Sean $z_1=r_1(\cos\theta_1+i\sin\theta_1)$ y $z_2=r_2(\cos\theta_2+i\sin\theta_2)$. Entonces
\[
z_1z_2=r_1r_2\bigl(\cos(\theta_1+\theta_2)+i\sin(\theta_1+\theta_2)\bigr),
\]
y, si $z_2\ne0$, $\dfrac{z_1}{z_2}=\dfrac{r_1}{r_2}\bigl(\cos(\theta_1-\theta_2)+i\sin(\theta_1-\theta_2)\bigr)$.
\end{teorema}

\begin{corolario}\label{cor:arg-producto}
Para $z,w\ne0$: $\abs{zw}=\abs{z}\abs{w}$ y $\arg(zw)=\arg z+\arg w$, en el sentido de que la suma de un argumento de $z$ y uno de $w$ es un argumento de $zw$. Para los argumentos principales, la igualdad solo vale salvo un múltiplo de $2\pi$.
\end{corolario}

\emph{Interpretación geométrica.} Fijado $w=\rho(\cos\alpha+i\sin\alpha)$, la aplicación $z\mapsto wz$ multiplica el módulo de cada punto por $\rho$ y suma $\alpha$ a su argumento: es la composición de un giro de ángulo $\alpha$ alrededor del origen y una homotecia de razón $\rho$. Esta lectura, que Needham sitúa en el centro de su exposición~\cite[cap.~1]{needham}, convierte en evidentes muchas identidades: por ejemplo, multiplicar por $i$ es girar un ángulo recto, y por eso $i^{2}=-1$ (dos giros rectos forman media vuelta).

\section{La exponencial imaginaria y la fórmula de Euler}

\begin{definicion}[Exponencial imaginaria]\label{def:euler}
Para $\theta\in\R$ se define
\[
e^{i\theta}=\cos\theta+i\sin\theta,
\]
y, para $z=x+iy\in\C$, $e^{z}=e^{x}e^{iy}=e^{x}(\cos y+i\sin y)$.
\end{definicion}

La notación no es arbitraria. Si se sustituye formalmente $x=i\theta$ en la serie $e^{x}=\sum_{n\ge0}x^{n}/n!$ y se agrupan los términos pares e impares usando $i^{2}=-1$, se obtienen exactamente las series de $\cos\theta$ y de $i\sin\theta$; las notas de Cambridge toman la serie como definición de $e^{z}$ y deducen la fórmula anterior~\cite[§1.2]{cambvm}. Las notas de 18.03 la justifican de otro modo: $f(\theta)=\cos\theta+i\sin\theta$ es la única función con $f'=if$ y $f(0)=1$, la misma ecuación que caracteriza a $e^{i\theta}$~\cite[§6.2]{mit1803}. Ambos argumentos requieren cálculo; aquí basta la definición y sus propiedades algebraicas.

\begin{proposicion}[Propiedades de $e^{i\theta}$]\label{prop:exp}
Para todos $\alpha,\beta,\theta\in\R$ y $n\in\Z$:
\begin{enumerate}
  \item $e^{i\alpha}e^{i\beta}=e^{i(\alpha+\beta)}$ y, más en general, $e^{z}e^{w}=e^{z+w}$ para $z,w\in\C$;
  \item $\abs{e^{i\theta}}=1$ y $\overline{e^{i\theta}}=e^{-i\theta}=1/e^{i\theta}$;
  \item $e^{i\theta}=1$ si y solo si $\theta\in2\pi\Z$; en particular, $e^{i\pi}=-1$ (identidad de Euler) y $e^{2\pi i}=1$;
  \item $\cos\theta=\dfrac{e^{i\theta}+e^{-i\theta}}{2}$ y $\sin\theta=\dfrac{e^{i\theta}-e^{-i\theta}}{2i}$.
\end{enumerate}
\end{proposicion}

Con esta notación, todo $z\ne0$ se escribe en \emph{forma exponencial} $z=re^{i\theta}$, con $r=\abs{z}$ y $\theta$ un argumento, y el \cref{teo:producto-polar} se resume en $r_1e^{i\theta_1}\cdot r_2e^{i\theta_2}=r_1r_2e^{i(\theta_1+\theta_2)}$. El conjugado es $\overline{re^{i\theta}}=re^{-i\theta}$.

\section{El teorema de De Moivre}

\begin{teorema}[De Moivre]\label{teo:demoivre}
Para todo $\theta\in\R$ y todo $n\in\Z$,
\[
(\cos\theta+i\sin\theta)^{n}=\cos n\theta+i\sin n\theta,\qquad\text{es decir,}\qquad (e^{i\theta})^{n}=e^{in\theta}.
\]
\end{teorema}

\begin{corolario}
Si $z=re^{i\theta}$, entonces $z^{n}=r^{n}e^{in\theta}$ para todo $n\in\Z$ (para $n<0$, con $z\ne0$).
\end{corolario}

\addtocounter{definicion}{1}%

\section{Raíces \texorpdfstring{$n$}{n}-ésimas y raíces de la unidad}

\begin{teorema}[Raíces $n$-ésimas]\label{teo:raices}
Sean $w=\rho e^{i\varphi}\ne0$ y $n\ge1$. La ecuación $z^{n}=w$ tiene exactamente $n$ soluciones distintas,
\[
z_k=\rho^{1/n}\,e^{i(\varphi+2\pi k)/n},\qquad k=0,1,\dots,n-1,
\]
donde $\rho^{1/n}$ es la raíz $n$-ésima real positiva de $\rho$.
\end{teorema}

Las $n$ raíces tienen el mismo módulo y argumentos separados $2\pi/n$: son los vértices de un polígono regular de $n$ lados centrado en el origen.

\begin{definicion}[Raíces de la unidad]
Las soluciones de $z^{n}=1$ son las \emph{raíces $n$-ésimas de la unidad}, $1,\omega,\omega^{2},\dots,\omega^{n-1}$ con $\omega=e^{2\pi i/n}$.
\end{definicion}

\begin{proposicion}\label{prop:suma-raices}
Si $n\ge2$, $\;1+\omega+\omega^{2}+\dots+\omega^{n-1}=0$.
\end{proposicion}

Geométricamente, la suma de los vectores que van del centro a los vértices de un polígono regular es nula: el conjunto es invariante por el giro de ángulo $2\pi/n$ (multiplicar la suma por $\omega$ la deja igual), y el único vector fijo por un giro no trivial es el cero.

\addtocounter{figure}{1}%

\section{Giros del plano como multiplicación por \texorpdfstring{$e^{i\alpha}$}{e\textasciicircum ia}}

\begin{proposicion}[Giros]\label{prop:giro-complejo}
Identifíquese $z=x+iy$ con el vector $(x,y)\T$. Para todo $\alpha\in\R$:
\begin{enumerate}
  \item la aplicación $z\mapsto e^{i\alpha}z$ corresponde a $(x,y)\T\mapsto R(\alpha)(x,y)\T$, donde $R(\alpha)$ es la matriz del \cref{cor:giros};
  \item conserva los módulos y los productos escalares: $\abs{e^{i\alpha}z}=\abs{z}$ y $\Real\bigl(e^{i\alpha}z\,\overline{e^{i\alpha}w}\bigr)=\Real(z\bar w)$;
  \item es biyectiva, con inversa $z\mapsto e^{-i\alpha}z$, y la composición de los giros de ángulos $\alpha$ y $\beta$ es el giro de ángulo $\alpha+\beta$.
\end{enumerate}
\end{proposicion}

La \cref{prop:giro-complejo}(iii) es la versión compleja del \cref{cor:giros}, y su demostración es más corta porque la regla de los exponentes ya contiene las fórmulas de la suma. Para girar alrededor de un punto $c$ distinto del origen se traslada, se gira y se deshace la traslación: $z\mapsto c+e^{i\alpha}(z-c)$.

\begin{corolario}[Productos escalares de vectores girados]\label{cor:rope}
Sean $z_1,z_2\in\C$, $\theta\in\R$ y $m,n\in\Z$. Entonces
\[
\Real\bigl(z_1e^{im\theta}\,\overline{z_2e^{in\theta}}\bigr)=\Real\bigl(z_1\bar z_2\,e^{i(m-n)\theta}\bigr),
\]
que solo depende de $m$ y $n$ a través de la diferencia $m-n$.
\end{corolario}

\addtocounter{definicion}{1}%

\addtocounter{definicion}{1}%

